\documentclass[10pt,a4paper]{amsart}
\usepackage[margin=19mm]{geometry}
\usepackage{amsmath,amssymb,mathtools}
\usepackage[scr=rsfs,cal=euler]{mathalfa}
\usepackage{xfrac}
\usepackage[unicode,hidelinks,bookmarks=false]{hyperref}
\newcommand{\ie}{i.e.\ }
\newcommand{\cf}{cf.\ }
\newcommand{\resp}{respectively\ }
\newcommand{\Subsection}{\subsection}
\providecommand{\nobreakdash}{\nobreak\hbox{-}}
\setlength{\emergencystretch}{2em}
\multlinegap0pt
\usepackage{soul}
\usepackage{tikz}
\usepackage{mathrsfs}
\usepackage{epstopdf}
\usepackage{placeins}
\usepackage{enumitem}
\usepackage{ulem}
\usetikzlibrary{calc,arrows.meta, positioning}
\numberwithin{equation}{section}
\allowdisplaybreaks
\newcommand{\esup}{\operatorname*{ess\,sup}}
\newcommand{\einf}{\operatorname*{ess\,inf}}
\newcommand{\eosc}{\operatorname*{ess\,osc}}
\newcommand{\rmnum}[1]{\romannumeral #1} % Lowercase roman number
\newcommand{\Rmnum}[1]{\uppercase\expandafter{\romannumeral#1}} % Uppercase roman number
\def\Xint#1{\mathchoice
{\XXint\displaystyle\textstyle{#1}}%
{\XXint\textstyle\scriptstyle{#1}}%
{\XXint\scriptstyle\scriptscriptstyle{#1}}%
{\XXint\scriptscriptstyle\scriptscriptstyle{#1}}%
\!\int}
\def\XXint#1#2#3{{\setbox0=\hbox{$#1{#2#3}{\int}$ }
\vcenter{\hbox{$#2#3$ }}\kern-.6\wd0}}
\def\ddashint{\Xint=}
\def\dashint{\Xint-}
\newtheorem{theorem}{Theorem}[section]
\newtheorem{proposition}[theorem]{Proposition}
\newtheorem{lemma}[theorem]{Lemma}
\newtheorem{corollary}[theorem]{Corollary}
\newtheorem{definition}[theorem]{Definition}
\newtheorem{remark}[theorem]{Remark}
\newtheorem{question}[theorem]{Question}
\newcommand{\Rd}{\color{red}}
\newcommand{\Bl}{\color{blue}}
\setcounter{equation}{0}
\begin{document}
\title[Elliptic Harnack inequality and Poincar'e inequality for p-e]{Elliptic Harnack inequality and Poincar\'e inequality for $p$-energies on metric measure spaces}
\author{Meng Yang}
\date{}
\maketitle
\noindent\emph{Source and licence.} Elliptic Harnack inequality and Poincar\'e inequality for $p$-energies on metric measure spaces. Version arXiv:2606.16405v1 (CC BY 4.0). This material is distributed under the Creative Commons Attribution 4.0 International licence, http://creativecommons.org/licenses/by/4.0/, as recorded in the arXiv OAI licence metadata for this exact version and shown on the abstract page https://arxiv.org/abs/2606.16405v1. Full rights notice: licenses/arxiv-2606.16405-CC-BY-4.0.txt.\par\smallskip
\noindent\textit{Editorial scope.} Counted unit: the original \texttt{proposition} environment \texttt{prop\_PI} at source lines 353--359 together with its complete proof at lines 361--397; the delivered document keeps source lines 125--398 so that every referenced label and every citation resolves inside it. Native editable source is the paper's own concatenated TeX; the AMS wrapper is editorial formatting only. No publisher class, upstream script or unrelated artwork is redistributed.
\begin{equation}\label{eq_equiv}
\text{EHI}+\text{cap}(\Psi)\Leftrightarrow\text{PI}(\Psi)+\text{CS}(\Psi),
\end{equation}
under the same structural assumptions.

As a consequence, the nonlinear theory now has the same structural form as the classical Dirichlet form theory at the elliptic level: the Harnack--capacity side and the Poincar\'e--cutoff side are equivalent. This gives a robust substitute, in the absence of heat kernels, for the classical equivalences known in the linear setting.

The proof combines the technique developed in our previous work \cite{Yan25d} with a very recent technique introduced in \cite{EK26}. In \cite{EK26}, the authors established, for the first time in the \emph{nonlinear} setting, the equivalence between the parabolic Harnack inequality $\text{PHI}$ and the conjunction of the volume doubling condition and the Poincar\'e inequality $\text{PI}$, on metric measure spaces endowed with an upper gradient structure. Certain arguments in \cite{EK26}, however, are sufficiently robust to be adapted to more general metric measure spaces, including fractals. Following \cite[Lemmas 2.5 and 7.5, Theorem 7.6]{EK26}, the proof of $\text{PI}(\Psi)$ can be reduced to establishing $\text{EHI}$ together with a version of the Faber--Krahn inequality $\text{FK}(\Psi)$, which they call the Poincar\'e-type inequality for functions with zero boundary values; see \cite[(7.1), (7.10)]{EK26}. In the setting of \cite{EK26}, both $\text{FK}$ and $\text{EHI}$ are obtained from $\text{PHI}$: $\text{FK}$ follows from \cite[Proposition 7.1]{EK26}, while $\text{PHI}$ directly implies $\text{EHI}$. In our setting, since $\text{EHI}$ is assumed, it remains to establish $\text{FK}(\Psi)$. Building on our previous result in \cite{Yan25d} on Wolff potential estimates under $\text{EHI}$, $\text{cap}(\Psi)$, and certain geometric and analytic assumptions, we first prove an estimate of the measure of a set in terms of its relative capacity. Combining this estimate with a standard Maz'ya decomposition technique, we obtain $\text{FK}(\Psi)$, and hence $\text{PI}(\Psi)$.

We now turn to the formal statement of our result. To begin with, let us fix the notation. Throughout this paper, $p\in(1,+\infty)$ is fixed. The letters $C,C_1,C_2,C_A, C_B$ will always refer to some positive constants and may change at each occurrence. The sign $\asymp$ means that the ratio of the two sides is bounded from above and below by positive constants. The sign $\lesssim$ ($\gtrsim$) means that the LHS is bounded by positive constant times the RHS from above (below). For $x,y\in \mathbb{R}$, denote $x\vee y=\max\{x,y\}$, $x\wedge y=\min\{x,y\}$, $x_+=\max\{x,0\}$, and $x_-=\max\{-x,0\}$. We use $\# A$ to denote the cardinality of a set $A$.

We say that a function $\Psi:[0,+\infty)\to[0,+\infty)$ is doubling if $\Psi$ is a homeomorphism, which implies that $\Psi$ is strictly increasing continuous and $\Psi(0)=0$, and there exists $C_\Psi>1$, called a doubling constant of $\Psi$, such that $\Psi(2r)\le C_\Psi\Psi(r)$ for any $r>0$. Throughout this paper, we always assume that $\Psi$ is a doubling function with a doubling constant $C_\Psi$, and there exist $\beta_*, \beta^*>0$ with $\beta_*\le\beta^*$ such that
$$\frac{1}{C_\Psi}\left(\frac{R}{r}\right)^{\beta_*}\le \frac{\Psi(R)}{\Psi(r)}\le C_{\Psi}\left(\frac{R}{r}\right)^{\beta^*}\text{ for any }r\le R.$$
Indeed, we can take $\beta^*=\log_2C_\Psi$.

Let $(X,d,m)$ be a complete unbounded metric measure space, that is, $(X,d)$ is a complete unbounded locally compact separable metric space and $m$ is a positive Radon measure on $X$ with full support. Throughout this paper, we always assume that all metric balls are relatively compact. For any $x\in X$ and any $r>0$, denote $B(x,r)=\{y\in X:d(x,y)<r\}$ and $V(x,r)=m(B(x,r))$. If $B=B(x,r)$, then denote $\delta B=B(x,\delta r)$ for any $\delta>0$. Let $\mathcal{B}(X)$ be the family of all Borel measurable subsets of $X$. Let $C(X)$ be the family of all continuous functions on $X$. Let $C_c(X)$ be the family of all continuous functions on $X$ with compact support. Denote $\dashint_A=\frac{1}{m(A)}\int_A$ and $u_A=\dashint_Au\mathrm{d} m$ for any measurable set $A$ with $m(A)\in(0,+\infty)$ and any function $u$ such that the integral $\int_Au\mathrm{d} m$ is well-defined.

We say that the chain condition \ref{eq_CC} holds if there exists $C_{cc}>0$ such that for any $x,y\in X$ and any integer $n\ge1$, there exists a sequence $\{x_k:0\le k\le n\}$ of points in $X$ with $x_0=x$ and $x_n=y$ such that
\begin{equation*}\label{eq_CC}\tag*{CC}
d(x_k,x_{k-1})\le C_{cc} \frac{d(x,y)}{n}\text{ for any }k=1,\ldots,n.
\end{equation*}
Throughout this paper, we always assume \ref{eq_CC}.

We say that the linearly locally connected condition \hypertarget{eq_LLC}{\text{LLC}} holds if there exists $A_{LLC}>1$ such that for any $x_0\in X$ and any $r>0$, for any $x,y\in B(x_0,r)\backslash B(x_0,r/2)$, there exists a connected compact set $K\subseteq B(x_0,A_{LLC}r)\backslash B(x_0,r/(2A_{LLC}))$ containing $x,y$. See also \cite[3.12]{HK98} or \cite[Page 234]{HKST15}.

We say that the volume doubling condition \ref{eq_VD} holds if there exists $C_{VD}>0$ such that
\begin{equation*}\label{eq_VD}\tag*{VD}
V(x,2r)\le C_{VD}V(x,r)\text{ for any }x\in X,r>0.
\end{equation*}
It follows directly from \ref{eq_VD} that, for any $x,y\in X$ and any $R,r>0$ with $r\le R$, we have
\begin{equation}\label{eq_VDa}
\frac{V(x,R)}{V(y,r)}\le C_{VD} \left(\frac{d(x,y)+R}{r}\right)^{\alpha},
\end{equation}
where $\alpha=\log_2C_{VD}$.

We say that $(\mathcal{E},\mathcal{F})$ is a $p$-energy on $(X,d,m)$ if $\mathcal{F}$ is a dense subspace of $L^p(X;m)$ and $\mathcal{E}:\mathcal{F}\to[0,+\infty)$ satisfies the following conditions.

\begin{enumerate}[label=(\arabic*)]
\item $\mathcal{E}^{1/p}$ is a semi-norm on $\mathcal{F}$, that is, for any $f,g\in\mathcal{F}$, $c\in\mathbb{R}$, we have $\mathcal{E}(f)\ge0$, $\mathcal{E}(cf)^{1/p}=|c|\mathcal{E}(f)^{1/p}$ and $\mathcal{E}(f+g)^{1/p}\le\mathcal{E}(f)^{1/p}+\mathcal{E}(g)^{1/p}$.
\item (Closed property) $(\mathcal{F},\mathcal{E}(\cdot)^{1/p}+\lVert {\cdot}\rVert_{L^p(X;m)})$ is a Banach space.
\item (Markovian property) For any $\varphi\in C(\mathbb{R})$ with $\varphi(0)=0$ and $|\varphi(t)-\varphi(s)|\le|t-s|$ for any $t,s\in\mathbb{R}$, for any $f\in\mathcal{F}$, we have $\varphi(f)\in\mathcal{F}$ and $\mathcal{E}(\varphi(f))\le\mathcal{E}(f)$.
\item (Regular property) $\mathcal{F}\cap C_c(X)$ is uniformly dense in $C_c(X)$ and $(\mathcal{E}(\cdot)^{1/p}+\lVert {\cdot}\rVert_{L^p(X;m)})$-dense in $\mathcal{F}$.
\item (Strongly local property) For any $f,g\in\mathcal{F}$ with compact support and $g$ constant in an open neighborhood of $\mathrm{supp}(f)$, we have $\mathcal{E}(f+g)=\mathcal{E}(f)+\mathcal{E}(g)$.
\item ($p$-Clarkson's inequality) For any $f,g\in\mathcal{F}$, we have
\begin{equation*}\label{eq_Cla}\tag*{Cla}
\begin{cases}
\mathcal{E}(f+g)+\mathcal{E}(f-g)\ge2 \left(\mathcal{E}(f)^{\frac{1}{p-1}}+\mathcal{E}(g)^{\frac{1}{p-1}}\right)^{p-1}&\text{if }p\in(1,2],\\
\mathcal{E}(f+g)+\mathcal{E}(f-g)\le2 \left(\mathcal{E}(f)^{\frac{1}{p-1}}+\mathcal{E}(g)^{\frac{1}{p-1}}\right)^{p-1}&\text{if }p\in[2,+\infty).\\
\end{cases}
\end{equation*}
\end{enumerate}
Moreover, we also always assume the following condition.
\begin{itemize}
\item ($\mathcal{F}\cap L^\infty(X;m)$ is an algebra) For any $f,g\in\mathcal{F}\cap L^\infty(X;m)$, we have $fg\in\mathcal{F}$ and
\begin{equation*}\label{eq_Alg}\tag*{Alg}
\mathcal{E}(fg)^{{1}/{p}}\le \lVert f\rVert_{L^\infty(X;m)}\mathcal{E}(g)^{1/p}+\lVert g\rVert_{L^\infty(X;m)}\mathcal{E}(f)^{1/p}.
\end{equation*}
\end{itemize}
Denote $\mathcal{E}_\lambda(\cdot)=\mathcal{E}(\cdot)+\lambda \lVert {\cdot}\rVert^p_{L^p(X;m)}$ for any $\lambda>0$. Indeed, a general condition called the generalized $p$-contraction property was introduced in \cite{KS24a}, which implies \ref{eq_Cla}, \ref{eq_Alg}, and holds on a large family of metric measure spaces.

We list some basic properties of $p$-energies as follows. For any $f,g\in\mathcal{F}$, the derivative
$$\mathcal{E}(f;g)=\frac{1}{p}\frac{\mathrm{d}}{\mathrm{d} t}\mathcal{E}(f+tg)|_{t=0}\in\mathbb{R}$$
exists, the map $\mathcal{E}(f;\cdot):\mathcal{F}\to\mathbb{R}$ is linear, $\mathcal{E}(f;f)=\mathcal{E}(f)$. Moreover, for any $f,g\in\mathcal{F}$, for any $a\in\mathbb{R}$, we have
\begin{equation}\label{eq_quasi_strict}
\mathbb{R}\ni t\mapsto\mathcal{E}(f+tg;g)\in\mathbb{R}\text{ is strictly increasing if and only if }\mathcal{E}(g)>0,
\end{equation}
$$\mathcal{E}(af;g)=\mathrm{sgn}(a)|a|^{p-1}\mathcal{E}(f;g),$$
$$|\mathcal{E}(f;g)|\le\mathcal{E}(f)^{(p-1)/p}\mathcal{E}(g)^{1/p}.$$
Moreover, for any $\lambda>0$, all of the above results remain valid with $\mathcal{E}$ replaced by $\mathcal{E}_\lambda$, and for any $f,g\in\mathcal{F}$, we have
$$\mathcal{E}_\lambda(f;g)=\mathcal{E}(f;g)+\lambda\int_X\mathrm{sgn}(f)|f|^{p-1}g\mathrm{d} m.$$
See \cite[Theorem 3.7, Corollary 3.25]{KS24a} for the proofs of these results.

By \cite[Theorem 2.4]{Sas26}, a $p$-energy $(\mathcal{E},\mathcal{F})$ corresponds to a (canonical) $p$-energy measure $\Gamma:\mathcal{F}\times\mathcal{B}(X)\to[0,+\infty)$, $(f,A)\mapsto\Gamma(f)(A)$ satisfying the following conditions.
\begin{enumerate}[label=(\alph*),ref=(\alph*)]
\item\label{item_meas1} For any $f\in\mathcal{F}$, $\Gamma(f)(\cdot)$ is a positive Radon measure on $X$ with $\Gamma(f)(X)=\mathcal{E}(f)$.
\item\label{item_meas2} For any $A\in\mathcal{B}(X)$, $\Gamma(\cdot)(A)^{1/p}$ is a semi-norm on $\mathcal{F}$.
\item\label{item_meas3} For any $f,g\in\mathcal{F}\cap C_c(X)$, $A\in\mathcal{B}(X)$, if $f-g$ is constant on $A$, then $\Gamma(f)(A)=\Gamma(g)(A)$.
\item\label{item_meas4} ($p$-Clarkson's inequality) For any $f,g\in\mathcal{F}$, for any $A\in\mathcal{B}(X)$, we have
\begin{equation*}
\begin{cases}
\Gamma(f+g)(A)+\Gamma(f-g)(A)\ge2 \left(\Gamma(f)(A)^{\frac{1}{p-1}}+\Gamma(g)(A)^{\frac{1}{p-1}}\right)^{p-1}&\text{if }p\in(1,2],\\
\Gamma(f+g)(A)+\Gamma(f-g)(A)\le2 \left(\Gamma(f)(A)^{\frac{1}{p-1}}+\Gamma(g)(A)^{\frac{1}{p-1}}\right)^{p-1}&\text{if }p\in[2,+\infty).\\
\end{cases}
\end{equation*}
\item\label{item_meas5} (Chain rule) For any $f,g\in\mathcal{F}\cap C_c(X)$, for any piecewise $C^1$ functions $\varphi,\psi:\mathbb{R}\to\mathbb{R}$ with $\varphi(0)=\psi(0)=0$, we have
$$\mathrm{d}\Gamma\left(\varphi(f);\psi(g)\right)=\mathrm{sgn}(\varphi'(f))|\varphi'(f)|^{p-1}\psi'(g)\mathrm{d}\Gamma(f;g),$$
where $\Gamma(f;g)$ is a signed measure given by $\Gamma(f;g)=\frac{1}{p}\frac{\mathrm{d}}{\mathrm{d}t}\Gamma(f+tg)|_{t=0}$.
\item\label{item_meas6} (Leibniz rule) For any $f,g,h\in \mathcal{F}\cap C_c(X)$, we have $\mathrm{d}\Gamma(f;gh)=g \mathrm{d}\Gamma(f;h)+h \mathrm{d}\Gamma(f;g)$.
\end{enumerate}
%Using the chain rule, we have the following condition.
%\begin{itemize}
%\item (Strong sub-additivity) For any $f,g\in\mathcal{F}$, we have $f\vee g, f\wedge g\in\mathcal{F}$ and
%\begin{equation*}\label{eq_SubAdd}\tag*{SubAdd}
%\mathcal{E}(f\vee g)+\mathcal{E}(f\wedge g)\le\mathcal{E}(f)+\mathcal{E}(g).
%\end{equation*}
%\end{itemize}

Let $A_1,A_2\in\mathcal{B}(X)$. We define the capacity between $A_1,A_2$ as
\begin{align*}
&\mathrm{cap}(A_1,A_2)=\inf\left\{\mathcal{E}(\varphi):\varphi\in\mathcal{F},
\begin{array}{l}
\varphi=1\text{ in an open neighborhood of }A_1,\\
\varphi=0\text{ in an open neighborhood of }A_2
\end{array}
\right\},
\end{align*}
here we use the convention that $\inf\emptyset=+\infty$.

We say that the two-sided capacity bounds \hypertarget{eq_cap}{$\text{cap}(\Psi)$} hold if both the capacity upper bound \ref{eq_ucap} and the capacity lower bound \ref{eq_lcap} hold as follows. There exist $C_{cap}>0$, $A_{cap}>1$ such that for any ball $B(x,r)$, we have
\begin{align*}
\mathrm{cap}\left(B(x,r),X\backslash B(x,A_{cap}r)\right)&\le C_{cap} \frac{V(x,r)}{\Psi(r)},\label{eq_ucap}\tag*{$\text{cap}(\Psi)_{\le}$}\\
\mathrm{cap}\left(B(x,r),X\backslash B(x,A_{cap}r)\right)&\ge \frac{1}{C_{cap}} \frac{V(x,r)}{\Psi(r)}.\label{eq_lcap}\tag*{$\text{cap}(\Psi)_{\ge}$}
\end{align*}

Let $U$ be an open subset of $X$. Let
$$\mathcal{F}(U)=\text{the }\mathcal{E}_1\text{-closure of }\mathcal{F}\cap C_c(U).$$
We say that $u\in\mathcal{F}$ is harmonic in $U$ if $\mathcal{E}(u;v)=0$ for any $v\in\mathcal{F}\cap C_c(U)$, denoted by $-\Delta_pu=0$ in $U$. We say that $u\in\mathcal{F}$ is superharmonic in $U$ (resp. subharmonic in $U$) if $\mathcal{E}(u;v)\ge0$ (resp. $\mathcal{E}(u;v)\le0$) for any non-negative $v\in\mathcal{F}\cap C_c(U)$, denoted by $-\Delta_pu\ge0$ in $U$ (resp. $-\Delta_pu\le0$ in $U$). Moreover, if $U$ is bounded, then by \cite[Lemma 6.2]{Yan25c}, we have
$$\mathcal{F}(U)=\left\{u\in\mathcal{F}:\widetilde{u}=0\text{ q.e. on }X\backslash U\right\}.$$
Assuming that $U$ is bounded, then the above equality for harmonic functions, as well as the corresponding inequalities for superharmonic and subharmonic functions, also hold for all $v\in\mathcal{F}(U)$.

We say that the bottom spectrum positivity condition \ref{eq_BSP} holds if for any bounded open subset $U$ of $X$, we have
\begin{equation*}\label{eq_BSP}\tag*{\text{BSP}}
\lambda_1(U)=\inf \left\{\frac{\mathcal{E}(u)}{\lVert u\rVert_{L^p(X;m)}^p}:u\in \mathcal{F}(U)\backslash \{0\}\right\}>0.
\end{equation*}

We say that the lower semi-continuous condition \hypertarget{eq_LSC}{\text{LSC}} holds if for any bounded open subset $U$ of $X$, for any $u\in \mathcal{F}$ which is bounded from below superharmonic in $U$, we have $u$ has a modification which is lower semi-continuous in $U$.

We say that the elliptic Harnack inequality \ref{eq_EHI} holds if there exist $C_H>0$, $A_H>1$ such that for any ball $B(x,r)$, for any $u\in\mathcal{F}$ which is non-negative harmonic in $B(x,A_Hr)$, we have
\begin{equation*}\label{eq_EHI}\tag*{\text{EHI}}
\esup_{B(x,r)}u\le C_H\einf_{B(x,r)}u.
\end{equation*}

We say that the annulus elliptic Harnack inequality \ref{eq_EHIann} holds if there exist $A_1$, $A_2$, $A_3>1$ with $A_1<A_2<A_3$ and $C_H>0$ such that the following holds: for any ball $B(x,r)$ and any $u\in \mathcal{F}$ which is non-negative superharmonic in $B(x,A_3r)$ and harmonic in $B(x,A_3r)\backslash \overline{B(x,r)}$, we have
\begin{equation*}\label{eq_EHIann}\tag*{$\text{EHI}_{\text{ann}}$}
\esup_{B(x,A_2r)\backslash B(x,A_1r)}u\le C_H\einf_{B(x,A_2r)\backslash B(x,A_1r)}u.
\end{equation*}
See also \cite[Theorem 3.1]{Bla01} and \cite[LEMMA 6.3]{GS05} for similar results. However, in addition to harmonicity in the annulus, we also require superharmonicity in the entire ball. By a standard chaining argument, as in these references, it is easy to see that
\begin{center}
\ref{eq_VD} + \hyperlink{eq_LLC}{\text{LLC}} + \ref{eq_EHI} $\Rightarrow$ \ref{eq_EHIann}.
\end{center}

We say that the Poincar\'e inequality \ref{eq_PI} holds if there exist $C_{PI}>0$, $A_{PI}\ge1$ such that for any ball $B(x,r)$, for any $f\in\mathcal{F}$, we have
\begin{equation*}\label{eq_PI}\tag*{PI($\Psi$)}
\int_{B(x,r)}\lvert f-f_{B(x,r)}\rvert^p\mathrm{d} m\le C_{PI}\Psi(r)\int_{{B(x,A_{PI}r)}}\mathrm{d}\Gamma(f).
\end{equation*}
Under \ref{eq_VD} and \ref{eq_PI}, both \ref{eq_BSP} and \ref{eq_lcap} hold; these follow from \cite[Lemma 4.1]{Yan25c} and \cite[Lemma 5.1 (a)]{BB04}, respectively.

The main result of this paper is as follows.

\begin{theorem}\label{thm_main}
Assume \ref{eq_VD}, \hyperlink{eq_LSC}{\text{LSC}}, \ref{eq_BSP}, \ref{eq_EHI}, \hyperlink{eq_cap}{$\text{cap}(\Psi)$}. If either
\begin{enumerate}[label=(\alph*),ref=(\alph*)]
\item\label{cond_main_LLC} \hyperlink{eq_LLC}{\text{LLC}}, or
\item\label{cond_main_EHIann} \ref{eq_EHIann},
\end{enumerate}
holds, then \ref{eq_PI} holds.
\end{theorem}


\begin{remark}
\begin{enumerate}[label=(\arabic*)]
\item Since \ref{eq_VD}, \hyperlink{eq_LLC}{\text{LLC}}, \ref{eq_EHI} imply \ref{eq_EHIann}, it suffices here to treat case \ref{cond_main_EHIann}.
\item For the case $p=2$, it is well-known in the theory of heat kernel estimates that, under \ref{eq_VD}, the conjunction of \ref{eq_EHI} and \hyperlink{eq_cap}{$\text{cap}(\Psi)$} already implies \ref{eq_PI}; see \cite[THEOREM 1.2]{GHL15}. However, for general $p>1$, due to the \emph{nonlinearity} and the limitation of our techniques, \hyperlink{eq_LSC}{\text{LSC}}, \ref{eq_BSP}, and either \hyperlink{eq_LLC}{\text{LLC}} or \ref{eq_EHIann} are additionally required.
\end{enumerate}
\end{remark}

Combining Theorem \ref{thm_main} with our earlier results in \cite{Yan25d}, we obtain several versions of (\ref{eq_equiv}). To state these results, we first introduce the cutoff Sobolev inequality.

Let $U,V$ be two open subsets of $X$ satisfying $U\subseteq\overline{U}\subseteq V$. We say that $\phi\in\mathcal{F}$ is a cutoff function for $U\subseteq V$ if $0\le\phi\le1$ in $X$, $\phi=1$ in an open neighborhood of $\overline{U}$ and $\mathrm{supp}(\phi)\subseteq V$, where $\mathrm{supp}(f)$ refers to the support of the measure of $|f|\mathrm{d} m$ for any given function $f$.

We say that the cutoff Sobolev inequality \ref{eq_CS} holds if there exist $C_{1},C_{2}>0$, $A_{S}>1$ such that for any ball $B(x,r)$, there exists a cutoff function $\phi\in\mathcal{F}$ for $B(x,r)\subseteq B(x,A_Sr)$ such that for any $f\in\mathcal{F}$, we have
\begin{equation*}\label{eq_CS}\tag*{CS($\Psi$)}
\int_{B(x,A_{S}r)}|\widetilde{f}|^p\mathrm{d}\Gamma(\phi)\le C_{1}\int_{B(x,A_{S}r)}\mathrm{d}\Gamma(f)+\frac{C_{2}}{\Psi(r)}\int_{B(x,A_{S}r)}|f|^p\mathrm{d} m,
\end{equation*}
where $\widetilde{f}$ is a quasi-continuous modification of $f$, such that $\widetilde{f}$ is uniquely determined $\Gamma(\phi)$-a.e. in $X$, see \cite[Section 8]{Yan25a} for more details.

%Under \ref{eq_VD}, by taking $f\equiv1$ in $B(x,A_Sr)$, it is easy to see that \ref{eq_CS} implies \ref{eq_ucap}.

We have the following equivalences.

\begin{corollary}\label{cor_EHIann}
Assume \ref{eq_VD}. The followings are equivalent.
\begin{enumerate}[label=(\arabic*),ref=(\arabic*)]
\item\label{item_EHIann1} \hyperlink{eq_LSC}{\text{LSC}}, \ref{eq_BSP}, \ref{eq_EHI}, \ref{eq_EHIann}, \hyperlink{eq_cap}{$\text{cap}(\Psi)$}.
\item\label{item_EHIann2} \ref{eq_PI}, \ref{eq_CS}.
\end{enumerate}
\end{corollary}

\begin{corollary}\label{cor_LLC}
Assume \ref{eq_VD}, \hyperlink{eq_LLC}{\text{LLC}}. The followings are equivalent.
\begin{enumerate}[label=(\roman*),ref=(\roman*)]
\item\label{item_LLC1} \hyperlink{eq_LSC}{\text{LSC}}, \ref{eq_BSP}, \ref{eq_EHI}, \hyperlink{eq_cap}{$\text{cap}(\Psi)$}.
\item\label{item_LLC2} \ref{eq_PI}, \ref{eq_CS}.
\end{enumerate}
\end{corollary}

\begin{proof}[Proofs of Corollaries \ref{cor_EHIann} and \ref{cor_LLC}]
The implications ``\ref{item_EHIann2}$\Rightarrow$\ref{item_EHIann1}" and ``\ref{item_LLC2}$\Rightarrow$\ref{item_LLC1}" follow directly from \cite[Theorem 2.3]{Yan25d}. For the converse implications, \ref{eq_PI} follows from Theorem \ref{thm_main}, while \ref{eq_CS} follows from \cite[Theorem 2.1]{Yan25d}.
\end{proof}


\section{Proof}\label{sec_proof}

We will frequently use the following elementary consequence of H\"older's inequality:
\begin{equation}\label{eq_ele}
\int_{B(x_0,R)}\lvert f-f_{B(x_0,R)}\rvert^p \mathrm{d}m\le 2^p\inf_{c\in \mathbb{R}}\int_{B(x_0,R)}\lvert f-c\rvert^p \mathrm{d}m,
\end{equation}
for any ball $B(x_0,R)$ and any function $f$ for which the integrals above are well-defined.

We divide the proof into two steps. In the first step, following the strategy of \cite{EK26}, we reduce the proof of \ref{eq_PI} to proving the following Faber--Krahn inequality.

We say that the Faber--Krahn inequality \ref{eq_FK} holds if there exists $C>0$ such that for any ball $B(x_0,R)$ and any $f\in \mathcal{F}(B(x_0,R))$, we have
\begin{equation*}\label{eq_FK}\tag*{$\text{FK}(\Psi)$}
\int_{B(x_0,R)}\lvert f\rvert^p \mathrm{d}m\le C\Psi(R)\int_{B(x_0,R)} \mathrm{d}\Gamma(f).
\end{equation*}
Obviously, \ref{eq_FK} implies \ref{eq_BSP}.

Let
\begin{align*}
&\mathcal{F}_{\mathrm{loc}}=\left\{u:
\begin{array}{l}
\text{for any relatively compact open set }U,\\
\text{there exists }u^\#\in \mathcal{F}\text{ such that }u=u^\#\text{ }m\text{-a.e. in }U
\end{array}
\right\}.
\end{align*}
For any $u\in \mathcal{F}_{\mathrm{loc}}$, let $\Gamma(u)|_U=\Gamma(u^\#)|_U$, where $u^\#, U$ are given as above, then $\Gamma(u)$ is a well-defined positive Radon measure on $X$, as follows from \ref{item_meas3} and \ref{item_meas6} in the definition of $\Gamma(u)$ for $u\in \mathcal{F}$, together with an argument similar to that in \cite[Corollary 3.2.1]{FOT11}.


\begin{proposition}\label{prop_PI}
Assume that there exist $C>0$, $A>1$, $\varepsilon>0$ with $\frac{1+\varepsilon}{2^{\frac{p-1}{p}}}<1$ such that for any ball $B(x_0,R)$ and any non-negative $f\in \mathcal{F}_{\mathrm{loc}}$, we have
\begin{equation}\label{eq_weakPI}
\left(\dashint_{B(x_0,R)}f^p \mathrm{d}m\right)^{1/p}\le C\left(\frac{\Psi(R)}{V(x_0,R)}\int_{B(x_0,AR)} \mathrm{d}\Gamma(f)\right)^{1/p}+(1+\varepsilon) \dashint_{B(x_0,R)}f \mathrm{d}m.
\end{equation}
Then \ref{eq_PI} holds.
\end{proposition}

\begin{proof}
For any ball $B(x_0,R)$ and any $f\in \mathcal{F}$, there exists $c\in \mathbb{R}$ such that
\begin{equation}\label{eq_medium}
\begin{aligned}
m \left(\left\{f\ge c\right\}\cap B(x_0,R)\right)&\ge \frac{1}{2}V(x_0,R),\\
m \left(\left\{f\le c\right\}\cap B(x_0,R)\right)&\ge \frac{1}{2}V(x_0,R).
\end{aligned}
\end{equation}
Applying (\ref{eq_weakPI}) to $(f-c)_+\in \mathcal{F}_{\mathrm{loc}}$, we have
\begin{align*}
&\left(\dashint_{B(x_0,R)}(f-c)_+^p \mathrm{d}m\right)^{1/p}\\
&\le C\left(\frac{\Psi(R)}{V(x_0,R)}\int_{B(x_0,AR)} \mathrm{d}\Gamma((f-c)_+)\right)^{1/p}+(1+\varepsilon) \dashint_{B(x_0,R)}(f-c)_+ \mathrm{d}m.
\end{align*}
By the Markovian property, we have
$$\int_{B(x_0,AR)} \mathrm{d}\Gamma((f-c)_+)\le \int_{B(x_0,AR)} \mathrm{d}\Gamma(f).$$
By H\"older's inequality and (\ref{eq_medium}), we have
\begin{align*}
&\dashint_{B(x_0,R)}(f-c)_+ \mathrm{d}m\\
&\le \left(\frac{m \left(\left\{f>c\right\}\cap B(x_0,R)\right)}{V(x_0,R)}\right)^{1-\frac{1}{p}} \left(\dashint_{B(x_0,R)}(f-c)_+^p \mathrm{d}m\right)^{\frac{1}{p}}\\
&\le \frac{1}{2^{\frac{p-1}{p}}}\left(\dashint_{B(x_0,R)}(f-c)_+^p \mathrm{d}m\right)^{\frac{1}{p}}.
\end{align*}
Hence
\begin{align*}
&\left(\dashint_{B(x_0,R)}(f-c)_+^p \mathrm{d}m\right)^{1/p}\\
&\le C\left(\frac{\Psi(R)}{V(x_0,R)}\int_{B(x_0,AR)} \mathrm{d}\Gamma(f)\right)^{1/p}+\frac{1+\varepsilon}{2^{\frac{p-1}{p}}}\left(\dashint_{B(x_0,R)}(f-c)_+^p \mathrm{d}m\right)^{{1}/{p}}.
\end{align*}
Let $\gamma=1-\frac{1+\varepsilon}{2^{\frac{p-1}{p}}}$, then by assumption, we have $\gamma\in(0,1)$, which gives
$$\int_{B(x_0,R)}(f-c)_+^p \mathrm{d}m\le \left(\frac{C}{\gamma}\right)^p{\Psi(R)}\int_{B(x_0,AR)} \mathrm{d}\Gamma(f).$$
Similarly,
$$\int_{B(x_0,R)}(f-c)_-^p \mathrm{d}m\le \left(\frac{C}{\gamma}\right)^p{\Psi(R)}\int_{B(x_0,AR)} \mathrm{d}\Gamma(f).$$
Therefore, by (\ref{eq_ele}),
\begin{align*}
&\int_{B(x_0,R)} \lvert f-f_{B(x_0,R)}\rvert^p \mathrm{d}m\le 2^p\int_{B(x_0,R)} \lvert f-c\rvert^p \mathrm{d}m\\
&=2^p\int_{B(x_0,R)} (f-c)_+^p \mathrm{d}m+2^p\int_{B(x_0,R)} (f-c)_-^p \mathrm{d}m\\
&\le 2^{p+1}\left(\frac{C}{\gamma}\right)^p{\Psi(R)}\int_{B(x_0,AR)} \mathrm{d}\Gamma(f).
\end{align*}
\end{proof}


\begin{thebibliography}{99}
\bibitem[BB04]{BB04} Martin~T. Barlow and Richard~F. Bass. \newblock Stability of parabolic {H}arnack inequalities. \newblock {\em Trans. Amer. Math. Soc.}, 356(4):1501--1533, 2004.
\bibitem[Bla01]{Bla01} S\'ebastien Blach\`ere. \newblock Harmonic functions on annuli of graphs. \newblock {\em Ann. Math. Blaise Pascal}, 8(2):47--59, 2001.
\bibitem[EK26]{EK26} Theo {Elenius} and Juha {Kinnunen}. \newblock {Doubling measures, Poincar{\'e} inequalities and parabolic Harnack inequalities for a doubly nonlinear equation}. \newblock {\em arXiv e-prints}, page arXiv:2605.27665, May 2026.
\bibitem[FOT11]{FOT11} Masatoshi Fukushima, Yoichi Oshima, and Masayoshi Takeda. \newblock {\em Dirichlet forms and symmetric {M}arkov processes}, volume~19 of {\em De Gruyter Studies in Mathematics}. \newblock Walter de Gruyter \& Co., Berlin, extended edition, 2011.
\bibitem[GHL15]{GHL15} Alexander Grigor'yan, Jiaxin Hu, and Ka-Sing Lau. \newblock Generalized capacity, {H}arnack inequality and heat kernels of {D}irichlet forms on metric measure spaces. \newblock {\em J. Math. Soc. Japan}, 67(4):1485--1549, 2015.
\bibitem[GS05]{GS05} Alexander Grigor'yan and Laurent Saloff-Coste. \newblock Stability results for {H}arnack inequalities. \newblock {\em Ann. Inst. Fourier (Grenoble)}, 55(3):825--890, 2005.
\bibitem[HK98]{HK98} Juha Heinonen and Pekka Koskela. \newblock Quasiconformal maps in metric spaces with controlled geometry. \newblock {\em Acta Math.}, 181(1):1--61, 1998.
\bibitem[HKST15]{HKST15} Juha Heinonen, Pekka Koskela, Nageswari Shanmugalingam, and Jeremy~T. Tyson. \newblock {\em Sobolev spaces on metric measure spaces}, volume~27 of {\em New Mathematical Monographs}. \newblock Cambridge University Press, Cambridge, 2015. \newblock An approach based on upper gradients.
\bibitem[KS24a]{KS24a} Naotaka {Kajino} and Ryosuke {Shimizu}. \newblock {Contraction properties and differentiability of $p$-energy forms with applications to nonlinear potential theory on self-similar sets}. \newblock {\em arXiv e-prints}, page arXiv:2404.13668v3, April 2024.
\bibitem[Sas26]{Sas26} K{\^o}hei {Sasaya}. \newblock {Construction of $p$-energy measures associated with strongly local $p$-energy forms}. \newblock {\em J. Funct. Anal.}, 291(1):Paper No. 111489, 2026.
\bibitem[Yan25a]{Yan25a} Meng {Yang}. \newblock {$p$-Poincar{\'e} inequalities and cutoff Sobolev inequalities on metric measure spaces}. \newblock {\em arXiv e-prints}, page arXiv:2504.09503v4, April 2025.
\bibitem[Yan25c]{Yan25c} Meng {Yang}. \newblock {On singularity of $p$-energy measures on metric measure spaces}. \newblock {\em arXiv e-prints}, page arXiv:2505.12468v2, May 2025.
\bibitem[Yan25d]{Yan25d} Meng {Yang}. \newblock {Energy inequalities for cutoff functions of $p$-energies on metric measure spaces}. \newblock {\em arXiv e-prints}, page arXiv:2507.08577v5, July 2025.
\end{thebibliography}
\end{document}
